\section{Matrices over \(\FF\bqty{x}\) and Normal Forms of Matrices}

In a previous example, we considered a vector space \(V\) over a field \(\FF\). Fix \(\alpha \in \End\pqty{V}\), we may construct a \(\FF\bqty{x}\) module
\[
    \function{\cdot}{\FF\bqty{x} \times V}{V}{\pqty{f, v}}{f\pqty{\alpha} \cdot v.}
\]

We call this module \(V_{\alpha}\).

If \(\dim V\) is finite, then \(V_{\alpha}\) is finitely generated as a \(\FF\bqty{x}\) module.

\begin{example}
    \label{ex:lower-diagonal-1}%
    Suppose \(V_{\alpha}\) is the module
    \[
        V_{\alpha} \isomorphic \flatfrac{\FF\bqty{x}}{\pqty{x^r}}.
    \]

    As a \(\FF\)-vector space, the right-hand side has basis
    \[
        1, x, x^2, \ldots, x^{r - 1}.
    \]  

    The endomorphism \(\functiondomain{\alpha}{V}{V}\) acts as multiplication by \(x\), i.e.,
    \[
        \function{\alpha}{\flatfrac{\FF\bqty{x}}{\pqty{x^r}}}{\flatfrac{\FF\bqty{x}}{\pqty{x^r}}}{f}{xf.}
    \]

    Take the endomorphism in the basis \(\Bqty{1, x, \ldots, x^{r - 1}}\) as a matrix
    \[
        \alpha = \pmqty{
            0 & 0 & \cdots & 0 & 0 & 0 \\
            1 & 0 & \cdots & 0 & 0 & 0 \\
            0 & 1 & \ddots & 0 & 0 & 0 \\
            \vdots & \vdots & \ddots & \ddots & \vdots & \vdots \\
            0 & 0 & \cdots & 1 & 0 & 0 \\
            0 & 0 & \cdots & 0 & 1 & 0
        }.
    \]
\end{example}

\begin{example}[Jordan Block]
    \label{ex:jordan-block}%
    Suppose now instead
    \[
        V_{\alpha} \isomorphic \flatfrac{\FF\bqty{x}}{\pqty{x - \lambda}^r}.
    \]

    Now we choose \(\functiondomain{\beta}{V}{V}\) such that \(\beta = \alpha - \lambda \identity\). Then there is a basis for \(V\), where
    \[
        \beta = \pmqty{
            0 & 0 & \cdots & 0 & 0 & 0 \\
            1 & 0 & \cdots & 0 & 0 & 0 \\
            0 & 1 & \ddots & 0 & 0 & 0 \\
            \vdots & \vdots & \ddots & \ddots & \vdots & \vdots \\
            0 & 0 & \cdots & 1 & 0 & 0 \\
            0 & 0 & \cdots & 0 & 1 & 0
        }.
    \]

    We note that in this basis, \(\alpha\) is a Jordan block of the form
    \[
        \alpha = \pmqty{
            \lambda & 0 & \cdots & 0 & 0 & 0 \\
            1 & \lambda & \cdots & 0 & 0 & 0 \\
            0 & 1 & \ddots & 0 & 0 & 0 \\
            \vdots & \vdots & \ddots & \ddots & \vdots & \vdots \\
            0 & 0 & \cdots & 1 & \lambda & 0 \\
            0 & 0 & \cdots & 0 & 1 & \lambda
        }.
    \]

    This is the \emph{Jordan block} of \emph{size} \(r\) and \emph{eigenvalue} \(\lambda\), and is denoted as \(\vb{J}_{r} \pqty{\lambda}\).
\end{example}

\begin{example}[Companion Matrix]
    \label{ex:companion-matrix}%
    Suppose now
    \[
        V_{\alpha} \isomorphic \flatfrac{\FF\bqty{x}}{p\pqty{x}}
    \]
    with \(\deg p = r\). A basis for the \(\FF\)-module on the right-hand side is given by
    \[
        1, x, x^2, \ldots, x^{r - 1}.
    \]

    Let
    \[
        p\pqty{x} = a_0 + a_1 x + \cdots + x^r,
    \]
    then the matrix for the map
    \[
        \function{\alpha}{\flatfrac{\FF\bqty{x}}{p\pqty{x}}}{\flatfrac{\FF\bqty{x}}{p\pqty{x}}}{f}{xf}
    \]
    is represented by, under this basis,
    \[
        \alpha = \pmqty{
            0 & 0 & \cdots & 0 & 0 & - a_0 \\
            1 & 0 & \cdots & 0 & 0 & - a_1 \\
            0 & 1 & \ddots & 0 & 0 & - a_2 \\
            \vdots & \vdots & \ddots & \ddots & \vdots & \vdots \\
            0 & 0 & \cdots & 1 & 0 & - a_{r - 2} \\
            0 & 0 & \cdots & 0 & 1 & - a_{r - 1}
        }.
    \]

    This is called the \emph{companion matrix} for \(p\pqty{x}\), and is denoted as \(\vb{C}\pqty{p}\).
\end{example}

\begin{theorem}[Rational Canonical Form]
    \label{thm:rational-canonical-form}%
    Let \(\alpha \in \End\pqty{V}\). Then
    \[
        V_{\alpha} \isomorphic \flatfrac{\FF\bqty{x}}{\pqty{f_1}} \oplus \cdots \oplus \flatfrac{\FF\bqty{x}}{\pqty{f_s}}
    \]
    with \(f_i \divides f_{i + 1}\).

    There exists a basis of \(V\), such that
    \[
        \alpha = \pmqty{\dmat{\vb{C} \pqty{f_1}, \vb{C} \pqty{f_2}, \ddots, \vb{C} \pqty{f_s}}} = \diag\pqty{\vb{C} \pqty{f_1}, \ldots, \vb{C} \pqty{f_s}}.
    \]
\end{theorem}

\begin{proof}
    We apply \autoref{thm:classification-finitely-generated-modules-euclidean-domains} (Structure Theorem for Finitely-Generated Modules over Euclidean Domains) for \(\FF\bqty{x}\) and \autoref{ex:companion-matrix} (Companion Matrix).
\end{proof}

\begin{remark}
    The \emph{minimal polynomial} of \(\alpha\) is \(f_s\), and the \emph{characteristic polynomial} of \(\alpha\) is the product of \(f_i\)s.
\end{remark}

\begin{lemma}[Fundamental Theorem of Algebra]
    \label{lem:fundamental-theorem-of-algebra}%
    The primes in \(\CC\bqty{x}\) are exactly \(\pqty{x - \alpha}\) for \(\alpha \in \CC\).
\end{lemma}

\begin{theorem}[Jordan Canonical Form]
    \label{thm:jordan-canonical-form}%
    Let \(\alpha \in \End\pqty{V}\), where \(V\) is a finite-dimensional vector space over \(\CC\). Then we have
    \[
        V_{\alpha} \isomorphic \flatfrac{\CC\bqty{x}}{\pqty{x - \lambda_1}^{\alpha_1}} \oplus \cdots \oplus \flatfrac{\CC\bqty{x}}{\pqty{x - \lambda_t}^{a_t}}.
    \]

    There exists a basis of \(V\), such that
    \[
        \alpha = \pmqty{\dmat{\vb{J}_{a_1} \pqty{\lambda_1}, \vb{J}_{a_2} \pqty{\lambda_2}, \ldots, \vb{J}_{a_t} \pqty{\lambda_t}}}
        = \diag\pqty{\vb{J}_{a_1} \pqty{\lambda_1}, \ldots, \vb{J}_{a_t} \pqty{\lambda_t}}.
    \]
\end{theorem}

\begin{proof}
    Apply \autoref{thm:primary-decomposition-classification} (Primary Decomposition of Structure Theorem) to \(V_{\alpha}\), use \autoref{lem:fundamental-theorem-of-algebra} (Fundamental Theorem of Algebra), and \autoref{ex:jordan-block} (Jordan Block).
\end{proof}

\begin{remark}
    We chose \(\CC\) here for the complete factorisation of a polynomial, since \(\CC\) is an algebraically closed field. However, Jordan Canonical Forms do exist for \emph{real} matrices as well -- but we can only factorise them into \emph{at most} quadratics, not necessarily linear terms. Things that look like \emph{rotation matrices} will come in as Jordan blocks for those, and they look slightly messier, but it is mathematically possible.
\end{remark}